% ===== PRÉAMBULE CANONIQUE — à coller VERBATIM en tête de CHAQUE .tex =====
\documentclass[11pt,a4paper]{article}
\usepackage[utf8]{inputenc}\usepackage[T1]{fontenc}\usepackage{lmodern}
\usepackage[french]{babel}
\usepackage{amsmath,amssymb,mathtools}\usepackage{icomma}
\usepackage{siunitx}\sisetup{locale=FR}  % \qty{}{}, \si{}, \ang{}
\usepackage{xcolor}\usepackage{colortbl}
\usepackage{tikz}\usepackage{pgfplots}\pgfplotsset{compat=1.18}
\usepackage[siunitx,europeanresistors]{circuitikz} % circuits (élec.)
\usepackage[most]{tcolorbox}\usepackage{enumitem}\usepackage{array,booktabs}
\usepackage[a4paper,margin=2.2cm]{geometry}
\usetikzlibrary{arrows.meta,calc,angles,quotes,positioning,patterns,%
  backgrounds,shapes.geometric,decorations.pathmorphing,decorations.markings,intersections}
\definecolor{primary}{RGB}{222,49,99}\definecolor{primarydark}{RGB}{150,22,55}
\definecolor{defcol}{RGB}{0,114,178}\definecolor{methcol}{RGB}{52,92,148}
\definecolor{excol}{RGB}{0,135,90}\definecolor{attncol}{RGB}{200,40,55}
\tcbset{boxbase/.style={enhanced,breakable,boxrule=0.9pt,arc=2.5pt,
  left=3mm,right=3mm,top=2.3mm,bottom=2.3mm,fonttitle=\bfseries,coltitle=white}}
% --- boîtes TUTORIEL (noms EXACTS, ne pas renommer) ---
\newtcolorbox{cle}[1][]{boxbase,colback=defcol!6,colframe=defcol,
  title={\textsc{À retenir}\ifx\relax#1\relax\relax\else\textmd{ : \itshape #1}\fi}}
\newtcolorbox{methode}[1][]{boxbase,colback=methcol!7,colframe=methcol,sharp corners,
  title={\textsc{Méthode}\ifx\relax#1\relax\relax\else\textmd{ --- \itshape #1}\fi}}
\newtcolorbox{exemplebox}[1][]{boxbase,colback=excol!5,colframe=excol,
  title={\textsc{Exemple}\ifx\relax#1\relax\relax\else\textmd{ : \itshape #1}\fi}}
\newtcolorbox{piege}[1][]{boxbase,colback=attncol!6,colframe=attncol,
  title={\textsc{Piège}\ifx\relax#1\relax\relax\else\textmd{ : \itshape #1}\fi}}
\newtcolorbox{remarque}[1][]{boxbase,colback=black!4,colframe=black!45,
  title={\textsc{Remarque}\ifx\relax#1\relax\relax\else\textmd{ : \itshape #1}\fi}}
% --- boîtes EXERCICE / CORRIGÉ (noms EXACTS, auto-comptées) ---
\newtcolorbox[auto counter]{exo}[1][]{boxbase,colback=primary!4,colframe=primary,
  title={\textsc{Exercice}~\thetcbcounter\ifx\relax#1\relax\relax\else\textmd{ --- \itshape #1}\fi}}
\newtcolorbox[auto counter]{corrige}[1][]{boxbase,colback=excol!4,colframe=excol,
  title={\textsc{Corrigé}~\thetcbcounter\ifx\relax#1\relax\relax\else\textmd{ --- \itshape #1}\fi}}
\newcommand{\vv}[1]{\vec{#1}}\newcommand{\ds}{\displaystyle}
\newcommand{\R}{\mathbb{R}}\newcommand{\N}{\mathbb{N}}\newcommand{\Z}{\mathbb{Z}}
% =========================================================================
\begin{document}
\sisetup{per-mode=symbol}

\begin{exo}[identifier la force de réaction]
Un livre est posé, immobile, sur une table elle-même posée sur le sol. Pour chaque force, donne sa
\textbf{réaction} au sens de la 3\textsuperscript{e} loi (nature, corps qui l'exerce, corps qui la subit).
\begin{enumerate}[label=\alph*)]
  \item La force de gravitation exercée par la Terre sur le livre.
  \item La réaction normale exercée par la table sur le livre.
  \item Compare l'intensité de la force gravitationnelle Terre~$\to$~livre et celle livre~$\to$~Terre.
\end{enumerate}
\end{exo}

\begin{exo}[peut-on frapper plus fort qu'on ne reçoit ?]
Une feuille de papier A4 est suspendue au plafond. Un élève tente de la frapper de façon à exercer sur
elle une force d'intensité \emph{supérieure} à celle que la feuille exerce sur lui.
\begin{enumerate}[label=\alph*)]
  \item Est-ce possible ? Justifie par la 3\textsuperscript{e} loi.
  \item La feuille étant très légère, comment expliquer qu'on ne puisse pas lui « donner » une grande force ?
\end{enumerate}
\end{exo}

\begin{exo}[Rida et William tirent une corde]
Rida et William tirent chacun à un bout d'une même corde ; William parvient à faire glisser Rida vers
lui. Réponds par \textbf{Vrai} ou \textbf{Faux} et justifie.
\begin{enumerate}[label=\Alph*.]
  \item La force que William exerce sur Rida (par la corde) est plus intense que celle que Rida exerce sur William.
  \item Aux deux bouts de la corde, les forces exercées ont le même sens.
  \item Rida exerce sur William une force de même intensité et de sens opposé.
  \item Rida glisse vers William parce que la force que William exerce sur le \emph{sol} est plus grande que celle de Rida sur le sol.
\end{enumerate}
\end{exo}

\begin{exo}[la monoroue électrique]
Une personne roule sur une monoroue électrique : la roue pousse le sol vers l'arrière, le sol pousse la
roue vers l'avant. Réponds par \textbf{Vrai} ou \textbf{Faux}.
\begin{enumerate}[label=\Alph*.]
  \item Les deux forces d'action--réaction ont la même intensité et des sens opposés.
  \item Il y a une accélération de la monoroue vers l'avant et une (du sol) vers l'arrière.
  \item Les intensités des deux forces action--réaction sont identiques.
  \item Les normes des deux accélérations (monoroue et sol) sont identiques.
\end{enumerate}
\end{exo}

\begin{exo}[le recul du canon, capstone]
Un canon de masse $m_{\text{canon}}=\qty{1200}{\kilogram}$ tire un boulet de $m_{\text{boulet}}=\qty{6}{\kilogram}$.
Dans le canon, la vitesse du boulet passe de $0$ à \qty{300}{\metre\per\second} en \qty{0,02}{\second}.
\begin{enumerate}[label=\alph*)]
  \item Calcule l'accélération du boulet, puis la force exercée sur lui.
  \item En déduis la force exercée sur le canon, puis son accélération de recul.
  \item Vérifie que $\dfrac{a_{\text{boulet}}}{a_{\text{canon}}}=\dfrac{m_{\text{canon}}}{m_{\text{boulet}}}$.
\end{enumerate}
\end{exo}
\end{document}
